\documentclass[10pt,a4paper]{amsart}
\usepackage[margin=19mm]{geometry}
\usepackage{amsmath,amssymb,mathtools}
\usepackage[scr=rsfs,cal=euler]{mathalfa}
\usepackage{xfrac}
\usepackage[unicode,hidelinks,bookmarks=false]{hyperref}
\newcommand{\ie}{i.e.\ }
\newcommand{\cf}{cf.\ }
\newcommand{\resp}{respectively\ }
\newcommand{\Subsection}{\subsection}
\providecommand{\nobreakdash}{\nobreak\hbox{-}}
\setlength{\emergencystretch}{2em}
\multlinegap0pt
\usepackage{bm}
\usepackage{bbm}
\usepackage{dsfont}
\usepackage{xspace}
\usepackage{cancel}
\usepackage{mathtools}
\usepackage{amsthm}
\makeatletter
\@ifundefined{theorem}{\newtheorem{theorem}{Theorem}}{}
\@ifundefined{lemma}{\newtheorem{lemma}[theorem]{Lemma}}{}
\makeatother
\makeatletter
\def\Ga{{\Gamma}}
\newcommand{\I}{{\rm i}}
\newcommand{\Otilde}{\widetilde{O}}
\renewcommand{\P}{\mathbb P}
\newcommand{\bC}{\mathbb C}
\newcommand{\bE}{\mathbb E}
\newcommand{\bP}{\mathbb P}
\newcommand{\bR}{\mathbb R}
\newcommand{\bZ}{\mathbb Z}
\newcommand{\cA}{\mathcal A}
\newcommand{\cC}{\mathcal C}
\newcommand{\cE}{\mathcal{E}}
\newcommand{\cG}{\mathcal G}
\newcommand{\cN}{\mathcal N}
\newcommand{\cR}{\mathcal R}
\expandafter\ifx\csname manuallemma\endcsname\relax
\newenvironment{manuallemma}[1]{%
	\renewcommand\themanuallemmainner{#1}%
	\manuallemmainner
}{\endmanuallemmainner}
\fi
\expandafter\ifx\csname manualquestion\endcsname\relax
\newenvironment{manualquestion}[1]{%
	\renewcommand\themanualquestioninner{#1}%
	\manualquestioninner
}{\endmanualquestioninner}
\fi
\expandafter\ifx\csname manualtheorem\endcsname\relax
\newenvironment{manualtheorem}[1]{%
	\renewcommand\themanualtheoreminner{#1}%
	\manualtheoreminner
}{\endmanualtheoreminner}
\fi
\newcommand{\one}{\boldsymbol{1}}
\makeatother
\title[Law of large numbers for the discriminant of random polynomi]{Law of large numbers for the discriminant of random polynomials}
\author{Marcus Michelen, Oren Yakir}
\date{Source: arXiv:2506.12206v1 (CC BY 4.0)}
\begin{document}
\maketitle

\noindent\emph{Source and licence.} Law of large numbers for the discriminant of random polynomials. Version arXiv:2506.12206v1 (CC BY 4.0), distributed under the Creative Commons Attribution 4.0 International licence, as recorded in the arXiv licence metadata for this exact version and shown on the abstract page https://arxiv.org/abs/2506.12206v1. Full rights notice: licenses/arxiv-2506.12206-CC-BY-4.0.txt.\par\smallskip

\noindent\textit{Editorial scope.} Counted unit: the Lemma whose source label is \texttt{lemma:comparison\_with\_a\_gaussian\_on\_the\_event\_C\_pm} in the article (source0.tex lines 1023--1040, source label \texttt{lemma:comparison\_with\_a\_gaussian\_on\_the\_event\_C\_pm}), delivered together with the article's own complete original proof of it (source0.tex lines 1041--1115, one \texttt{proof} environment of 75 lines). No mathematics is written, repaired or re-derived: statement and proof are the article's own text line for line. Required context retained uncounted: lemma:Salem-Zygmund (lines 521--527); eq:G-def (lines 523--525); lemma:berry\_esseen\_in\_annulus (lines 888--901). Every definition, display and label of those lines is retained unchanged. Omitted from this package and not redistributed: 1--520, 526--887, 902--1022, 1116--1702. Nothing inside a retained excerpt is omitted or altered: every retained body is delivered exactly as the article's lines are written, so no omission receipt of any kind accompanies this package. Citations: the retained excerpts carry no citation command at all, so no bibliography entry of the article is transcribed and no reference is redistributed. Reference adaptation: none was needed and none was made; every reference inside the delivered unit resolves inside it, so no unresolved reference or citation is produced. Printed slips kept literal: no printed word, symbol, hypothesis or number of the counted statement or of its proof was altered. Layout and class: the article's own class is \texttt{amsart} with packages including amssymb, amsfonts, amsthm, amsmath, booktabs, graphicx, verbatim, centernot, xcolor, multicol, enumitem, geometry, fancyhdr, hyperref; the wrapper is the lane's pinned \texttt{amsart} editorial wrapper at 10pt on A4, which restates the theorem-like environments the retained text needs and adds the article's own macro definitions that the retained text needs (\texttt{\textbackslash Ga}, \texttt{\textbackslash I}, \texttt{\textbackslash Otilde}, \texttt{\textbackslash P}, \texttt{\textbackslash bC}, \texttt{\textbackslash bE}, \texttt{\textbackslash bP}, \texttt{\textbackslash bR}, \texttt{\textbackslash bZ}, \texttt{\textbackslash cA}, \texttt{\textbackslash cC}, \texttt{\textbackslash cE}, \texttt{\textbackslash cG}, \texttt{\textbackslash cN}, \texttt{\textbackslash cR}, \texttt{\textbackslash one}), with extra wrapper packages (bm, bbm, dsfont, xspace, cancel, mathtools). The delivered numbering is the unit's own. Assets: no retained span contains a figure environment, a graphics-inclusion command, a PDF-page inclusion, a table or a TikZ picture; no image file is copied into this package, the package declares no asset dependency and no image file is redistributed with it. Licence: version arXiv:2506.12206v1 of this article is distributed under the Creative Commons Attribution 4.0 International licence, as recorded in the arXiv licence metadata for this exact version and shown on the abstract page https://arxiv.org/abs/2506.12206v1. A full rights notice is delivered at licenses/arxiv-2506.12206-CC-BY-4.0.txt. Characters: every retained excerpt is pure ASCII, so the delivered engine, XeLaTeX with its default Latin Modern text font, reports no missing character.
    \begin{lemma} \label{lemma:Salem-Zygmund}
    For the event
     \begin{equation} \label{eq:G-def}
    \mathcal{G} = \bigcap_{\ell=0}^{2} \Big\{ \max_{\theta\in[0,2\pi]} \big|f_n^{(\ell)}\big((1+1/n) e^{\I \theta}\big)\big| \le n^{\ell + 1/2} \log n\Big\}  \, ,
    \end{equation}
    there exists $c>0$ so that $\P(\cG) \ge 1 -  e^{-c\log^2 n}$.
    \end{lemma}

    \begin{lemma}
        \label{lemma:berry_esseen_in_annulus}
        Let $z\in \cA\setminus \mathcal{S}$, then
        \begin{equation*}
        \sup_{K \subset \bC^2}\Big| \P\Big( \big(f_n(z),f_n'(z)\big) \in K \Big) - \P_g\Big( \big(g_n(z),g_n'(z)\big) \in K \Big) \Big| \leq 
        \Otilde(n^{-1/2})
        \end{equation*}
        where the supremum is over all convex sets $K\subset \bC^2$. Furthermore, for $z,w\in \cA\setminus \mathcal{S}$ such that $|z-w|\ge n^{-1/2}$, we have
        \begin{equation*}
        \sup_{K \subset \bC^4}\Big| \P\Big( \big(f_n(z),f_n'(z),f_n(w),f_n'(w)\big) \in K \Big) - \P_g\Big( \big(g_n(z),g_n'(z),g_n(w),g_n'(w)\big) \in K \Big) \Big| \leq 
        \Otilde(n^{-1/2})
        \end{equation*}
        where here the supremum is over all convex sets $K\subset \bC^4$. 
        \end{lemma}

    \begin{lemma}
        \label{lemma:comparison_with_a_gaussian_on_the_event_C_pm}
        For $z\in \cN$ we have
        \begin{equation}
        \label{eq:lemma_comparison_with_a_gaussian_on_the_event_C_pm_first}
            \Big| \bP(\cC_z^+) - \bP_g(\cC_z^+) \Big| \le n^{-1/20} \, .
        \end{equation}
        Furthermore, for both $\bullet\in\{-,+\}$ we have
        \begin{equation}
        \label{eq:lemma_comparison_with_a_gaussian_on_the_event_C_pm_second}
            \bigg| \bE\Big[ \log_\bullet \Big|\frac{f_n^\prime(z)}{n^{3/2}}\Big| \cdot \one_{\cC_z^+} \Big]  - \bE_g\Big[ \log_\bullet \Big|\frac{g_n^\prime(z)}{n^{3/2}}\Big| \cdot \one_{\cC_z^+} \Big] \, \bigg|\le n^{-1/20} \, .
        \end{equation}
        Finally, for $z,w\in \cN$ such that $|z-w|\ge n^{-1/2}$ and for $\bullet\in\{-,+\}$ we have
        \begin{equation}
        \label{eq:lemma_comparison_with_a_gaussian_on_the_event_C_pm_third}
            \bigg| \bE\Big[ \log_\bullet \Big|\frac{f_n^\prime(z)}{n^{3/2}}\Big| \cdot \log_\bullet \Big|\frac{f_n^\prime(w)}{n^{3/2}}\Big| \cdot \one_{\cC_z^+\cap \cC_w^{+}} \Big]  - \bE_g\Big[ \log_\bullet \Big|\frac{g_n^\prime(z)}{n^{3/2}}\Big| \cdot \log_\bullet \Big|\frac{g_n^\prime(w)}{n^{3/2}}\Big| \cdot \one_{\cC_z^+\cap \cC_w^{+}} \Big] \, \bigg|\le n^{-1/20} \, .
        \end{equation}
    \end{lemma}

    \begin{proof}
        We will only prove~\eqref{eq:lemma_comparison_with_a_gaussian_on_the_event_C_pm_third} as both~\eqref{eq:lemma_comparison_with_a_gaussian_on_the_event_C_pm_first} and~\eqref{eq:lemma_comparison_with_a_gaussian_on_the_event_C_pm_second} have very similar (and simpler) proofs. We will also just consider the case $\bullet = +$, as the complementary case $\bullet = -$ is identical. We can intersect both expectations with the event $\cG$ given by~\eqref{eq:G-def}, since Lemma~\ref{lemma:Salem-Zygmund} guarantees that
        \[
        \bE\Big[ \log_\bullet \Big|\frac{f_n^\prime(z)}{n^{3/2}}\Big| \cdot \log_\bullet \Big|\frac{f_n^\prime(w)}{n^{3/2}}\Big| \cdot \one_{\cC_z^+\cap \cC_w^{+}} \cdot \one_{\cG^c} \Big] \le \widetilde{O}\big(\bP(\cG^c)^{1/2}\big) \lesssim e^{-c\log^2 n} \, ,
        \]
        which we can safely absorb in the error term. For $z\in \cN$, we set 
        \begin{equation}
        \label{eq:def_of_E_z}
            \mathcal{E}_z = \Big\{ (X,Y) \in \bC^2 \, : \,  \frac{X}{nY} \in \cR_z^+ -z \, , \, \frac1{\log^9 n} \le |Y| \le \log^2 n \, , \, |X|\le \log^2 n \, \Big\} \, .
        \end{equation}
        That is, $\cE_z$ is the subset of $\bC^2$ which corresponds to the event $\cC_z^+\cap\cG$, with 
        \[
        X = \frac{f_n(z)}{\sqrt{n}} \qquad \text{and} \qquad Y = \frac{f_n^\prime(z)}{n^{3/2}} \, .
        \]
        Recalling that $\cR_z^+$ is a polar rectangle around $z$ with side lengths $\widetilde{O}(n^{-1-\beta})$, we get that $m(\cR_z^+) = \widetilde{O}(n^{-2-2\beta})$ and hence
        \begin{equation}
            \label{eq:bound_on_measure_of_E_z}
            m(\cE_z) \le \widetilde{O}\big(n^2 m(\cR_z^+) \big) = \widetilde{O}(n^{-2\beta}) \, .
        \end{equation}We further denote by $\mu$ the law of
        \[
        \Big(\frac{f_n(z)}{\sqrt{n}},\frac{f_n^\prime(z)}{n^{3/2}},\frac{f_n(w)}{\sqrt{n}},\frac{f_n^\prime(w)}{n^{3/2}}\Big) \qquad \text{on } \ \bC^4\, ,
        \]
        and we get that 
        \begin{multline}
        \label{eq:expectaion_of_logs_in_complex_coordinates}
            \bE\Big[ \log_+ \Big|\frac{f_n^\prime(z)}{n^{3/2}}\Big| \cdot \log_+ \Big|\frac{f_n^\prime(w)}{n^{3/2}}\Big| \cdot \one_{\cC_z^+\cap \cC_w^{+}} \cdot \one_{\cG} \Big] \\  = \iint_{\substack{  (X_1,Y_1)\in \cE_z \\  (X_2,Y_2)\in \cE_w }} \log_+|Y_1| \cdot \log_+|Y_2| \, {\rm d}\mu(X_1,Y_1,X_2,Y_2) \, .
        \end{multline}
        We will also denote by $\mu_g$ the corresponding (Gaussian) law of 
        \[
        \Big(\frac{g_n(z)}{\sqrt{n}},\frac{g_n^\prime(z)}{n^{3/2}},\frac{g_n(w)}{\sqrt{n}},\frac{g_n^\prime(w)}{n^{3/2}}\Big) \qquad \text{on }  \ \bC^4 \, . 
        \]
        In view of~\eqref{eq:expectaion_of_logs_in_complex_coordinates}, we want to show that 
        \begin{multline}
        \label{eq:proof_of_lemma_comparison_for_C_what_we_want_in_complex_coordinates}
            \iint_{\substack{  (X_1,Y_1)\in \cE_z \\  (X_2,Y_2)\in \cE_w }} \log_+|Y_1| \cdot \log_+|Y_2| \, {\rm d}\mu(X_1,Y_1,X_2,Y_2) \\ = \iint_{\substack{  (X_1,Y_1)\in \cE_z \\  (X_2,Y_2)\in \cE_w }} \log_+|Y_1| \cdot \log_+|Y_2| \, {\rm d}\mu_g(X_1,Y_1,X_2,Y_2) + O(n^{-1/20})
        \end{multline}
        which we do in what follows. Let $\kappa=\beta>0$ be a small parameter, and consider the tiling of $\bC^2$ parametrized by the lattice $(n^{-\kappa} 
 \, \bZ)^4$, namely
        \begin{equation*}
            \mathcal{Q} = \Big\{ p + \big[0,n^{-\kappa}]^{4} \, : \, p\in (n^{-\kappa} \, \bZ)^4 \Big\} \, .
        \end{equation*}
        Then all cubes in $\mathcal{Q}$ have mutually disjoint interiors and their union is $\bR^4 \simeq \bC^2$.  For $z\in \cN$ and $\cE_z$ given by~\eqref{eq:def_of_E_z} we set 
        \begin{equation*}
            \Gamma_z = \big\{ Q\in \mathcal{Q} \, : \, Q\subset \cE_z\big\} \, , \quad 
            \Gamma_z^{b} = \big\{ Q\in \mathcal{Q} \, : \, Q\cap \cE_z \not = \emptyset \, , \ Q^c\cap \cE_z \not=\emptyset \big\}  \, , \quad   \overline \Ga_z = \Ga_z \cup \Ga_z^b \, , 
        \end{equation*}
        and note that since $\cE_z$ is a Lipschitz domain,~\eqref{eq:bound_on_measure_of_E_z} implies that 
        \begin{equation}
        \label{eq:bound_on_caridinality_of_Gamma}
            |\Ga_z| \le \widetilde{O}(n^{4\kappa - 2\beta}) \quad \text{and} \quad |\Ga_z^b| \le  \widetilde{O}(n^{3\kappa-2\beta}) \, .
        \end{equation}                
        To prove~\eqref{eq:proof_of_lemma_comparison_for_C_what_we_want_in_complex_coordinates}, we need to compare between the probability measures $\mu$ and $\mu_g$. Indeed, for any pair of cubes $Q_1,Q_2\in \mathcal{Q}$, Lemma~\ref{lemma:berry_esseen_in_annulus} implies that
        \begin{equation}
        \label{eq:comparison_of_measures_on_lattice_cubes_complex_coordinates}
            \mu\big(Q_1\times Q_2\big) = \mu_g\big(Q_1\times Q_2\big) + \widetilde{O}(n^{-1/2}) \, .
        \end{equation}
        Furthermore, since $\mu_g$ is absolutely continuous with respect to Lebesgue measure in $\bC^4$ with a bounded (Gaussian) density, we also have
        \begin{equation}
        \label{eq:bound_on_gaussian_meausure_on_lattice_cubes}
           \mu_g\big(Q_1\times Q_2\big) = \widetilde{O}(n^{-8\kappa}) \, , 
        \end{equation}
        so by taking $\kappa>0$ sufficiently small~\eqref{eq:bound_on_gaussian_meausure_on_lattice_cubes} also applies for $\mu$ in place of $\mu_g$. Since $\log_+|\cdot|$ has uniformly bounded variation in $\bC$, an immediate consequence of~\eqref{eq:comparison_of_measures_on_lattice_cubes_complex_coordinates} is that 
        \begin{equation}
        \label{eq:comparison_of_measures_with_logarithmic_integral}
            \iint_{Q_1\times Q_2} \log_+|Y_1| \cdot \log_+|Y_2| \, {\rm d}\mu \\ = \iint_{Q_1\times Q_2} \log_+|Y_1| \cdot \log_+|Y_2| \, {\rm d}\mu_g + \widetilde{O}(n^{-8\kappa})
        \end{equation}
        for any $Q_1\in \overline{\Ga}_z$ and $Q_2\in \overline{\Ga}_w$. We get that
        \begin{align*}
            \iint_{\substack{  (X_1,Y_1)\in \cE_z \\  (X_2,Y_2)\in \cE_w }} & \log_+|Y_1| \cdot \log_+|Y_2| \, {\rm d}\mu(X_1,Y_1,X_2,Y_2) \\ & \le \sum_{Q_1\in \overline{\Ga}_z} \sum_{Q_2\in \overline{\Ga}_z} \iint_{Q_1\times Q_2} \log_+|Y_1| \cdot \log_+|Y_2| \, {\rm d}\mu \\ (\text{by}~\eqref{eq:comparison_of_measures_with_logarithmic_integral}) &\le \sum_{Q_1\in \overline{\Ga}_z} \sum_{Q_2\in \overline{\Ga}_z} \iint_{Q_1\times Q_2} \log_+|Y_1| \cdot \log_+|Y_2| \, {\rm d}\mu_g + \widetilde{O}\big(|\overline{\Ga}_z|^2 n^{-8\kappa}\big) 
            \\ (\text{by}~\eqref{eq:bound_on_caridinality_of_Gamma}) & \le \sum_{Q_1\in \Ga_z} \sum_{Q_2\in \Ga_z} \iint_{Q_1\times Q_2} \log_+|Y_1| \cdot \log_+|Y_2| \, {\rm d}\mu_g + \widetilde{O}\big(n^{-4\beta} + n^{-2\kappa -2\beta} \big) 
            \\ & \le 
            \iint_{\substack{  (X_1,Y_1)\in \cE_z \\  (X_2,Y_2)\in \cE_w }} \log_+|Y_1| \cdot \log_+|Y_2| \, {\rm d}\mu_g(X_1,Y_1,X_2,Y_2)  + \widetilde{O}(n^{-4\beta}) \, , 
        \end{align*}
        which proves the upper bound in~\eqref{eq:proof_of_lemma_comparison_for_C_what_we_want_in_complex_coordinates}, for $\beta>0$ sufficiently small. The lower bound also follows by switching between $\mu$ and $\mu_g$ is the above derivation, and altogether we have proved~\eqref{eq:proof_of_lemma_comparison_for_C_what_we_want_in_complex_coordinates}, which concludes the proof of the lemma.
    \end{proof}

\end{document}
